% Standalone insertion for gamma_boundary_revision.tex.
% Required in the preamble if these commands are not already present:
% \providecommand{\Si}{\operatorname{Si}}
% \providecommand{\Cin}{\operatorname{Cin}}
% \providecommand{\Shi}{\operatorname{Shi}}
% \providecommand{\Cinh}{\operatorname{Cinh}}
% Suggested bibliography keys:
%   AdamczewskiFaverjon2025 = B. Adamczewski and C. Faverjon,
%     Algebraic independence measures for values of E-functions and
%     M-functions, arXiv:2502.09999 (2025).
%   Delaygue2025 = E. Delaygue, A Lindemann--Weierstrass theorem for
%     E-functions, arXiv:2210.12046v2 (2025).

\providecommand{\Si}{\operatorname{Si}}
\providecommand{\Cin}{\operatorname{Cin}}
\providecommand{\Shi}{\operatorname{Shi}}
\providecommand{\Cinh}{\operatorname{Cinh}}
\providecommand{\Eexp}{\mathcal E_{\mathrm{exp}}}

\section{Pure exponential base change and special-function periods}
\label{sec:efunction-consequences}

The toric relation theorem has a useful global form.  Define the pure
exponential field
\begin{equation}\label{eq:pure-exponential-field}
 \Eexp:=\Qbar\bigl(e^\alpha:\alpha\in\Qbar\bigr)\subset\C.
\end{equation}
Although this is an infinitely generated field, every coefficient of a
polynomial relation belongs to a subfield generated by only finitely many
algebraic exponentials.  This elementary observation removes the toric
coordinates altogether.

\begin{theorem}[Pure exponential base change]\label{thm:pure-exp-base}
The family
\[
 \bigl\{\Ein(\lambda):\lambda\in\Qbar^\times\bigr\}
\]
is algebraically independent over \(\Eexp\), in the finite-subset sense.
Equivalently, for any distinct
\(\lambda_1,\ldots,\lambda_s\in\Qbar^\times\),
\[
 \operatorname{trdeg}_{\Eexp}
 \Eexp\bigl(\Ein(\lambda_1),\ldots,\Ein(\lambda_s)\bigr)=s.
\]
\end{theorem}

\begin{proof}
Suppose that a polynomial relation exists.  Its finitely many coefficients
belong to
\(\Qbar(e^\omega:\omega\in\Omega)\) for some finite
\(\Omega\subset\Qbar\).  Clear coefficient denominators and apply
Theorem~\ref{thm:master} to
\[
 \Sigma=\{\lambda_1,\ldots,\lambda_s\}
        \cup(\Omega\setminus\{0\}).
\]
That theorem makes all \(\Ein(\sigma)\), \(\sigma\in\Sigma\),
algebraically independent over
\(\Qbar(e^\sigma:\sigma\in\Sigma)\).  In particular the selected
\(\Ein(\lambda_j)\) are independent over the coefficient field, a
contradiction.
\end{proof}

We shall also use the following elementary linear-image consequence.  It is
recorded because it identifies full relation ideals, rather than only their
transcendence degrees.

\begin{lemma}[Linear images of the \(\Ein\)-coordinates]
\label{lem:linear-image-ein}
Let \(Y=(\Ein(\lambda_1),\ldots,\Ein(\lambda_s))^t\), where the
\(\lambda_j\) are distinct nonzero algebraic numbers, and let
\(B\in M_{r,s}(\Qbar)\).  Put \(W=BY\).  The kernel of
\[
 \Eexp[T_1,\ldots,T_r]\longrightarrow\C,
 \qquad T_j\longmapsto W_j,
\]
is the linear ideal
\[
 \left(\sum_{j=1}^r c_jT_j:
       c\in\Qbar^r,\ c^tB=0\right).
\]
In particular, the \(W_j\) are algebraically independent over \(\Eexp\)
when \(B\) has full row rank.
\end{lemma}

\begin{proof}
By Theorem~\ref{thm:pure-exp-base}, the entries of \(Y\) are independent
indeterminates over \(\Eexp\).  After invertible row and column operations
over \(\Qbar\), the map \(Y\mapsto BY\) is a coordinate projection of rank
\(\operatorname{rank}B\).  Its polynomial kernel is therefore exactly the
displayed linear ideal.
\end{proof}

\subsection{A graphic matroid of definite periods}

Let \(G=(V,E)\) be a finite oriented graph, let \(c(G)\) be its number of
connected components, and attach pairwise distinct labels
\(\lambda_v\in\Qbar^\times\) to its vertices.  For an edge \(e\), write
\(s(e)\) and \(t(e)\) for its source and target and set
\begin{equation}\label{eq:graph-period}
 P_e:=\int_{\lambda_{s(e)}}^{\lambda_{t(e)}}
             \frac{1-e^{-z}}{z}\,dz
     =\Ein(\lambda_{t(e)})-\Ein(\lambda_{s(e)}).
\end{equation}
The integrand is entire after filling its removable singularity at zero, so
the integral is path independent.

Let \(\partial:\Qbar^E\to\Qbar^V\) be the oriented incidence map,
\(\partial e=t(e)-s(e)\).  Its kernel is the cycle space of \(G\).

\begin{theorem}[Graphic-period relation ideal]\label{thm:graphic-period}
The kernel of
\[
 \Eexp[X_e:e\in E]\longrightarrow\C,
 \qquad X_e\longmapsto P_e,
\]
is
\begin{equation}\label{eq:cycle-ideal}
 \left(\sum_{e\in E}c_eX_e:c\in\ker\partial\right).
\end{equation}
Thus every algebraic relation among the periods \(P_e\) is generated by the
signed cycle sums.  In particular,
\[
 \operatorname{trdeg}_{\Eexp}\Eexp(P_e:e\in E)
   =|V|-c(G),
\]
and the periods attached to any spanning forest are algebraically
independent over \(\Eexp\).
\end{theorem}

\begin{proof}
Put \(Y_v=\Ein(\lambda_v)\).  In vector notation,
\((P_e)_{e\in E}=\partial^t(Y_v)_{v\in V}\).  Theorem
\ref{thm:pure-exp-base} makes the \(Y_v\) algebraically independent over
\(\Eexp\), so Lemma~\ref{lem:linear-image-ein} identifies the kernel with
the linear relations among the rows of \(\partial^t\), namely
\(\ker\partial\).  The incidence matrix has rank \(|V|-c(G)\), and its
cycle space is spanned by signed circuit vectors.  This proves every claim.
\end{proof}

\subsection{The fourth-root orbit}

Define the four entire integral functions
\begin{equation}\label{eq:four-entire-integrals}
\begin{aligned}
 \Si(z)&=\int_0^z\frac{\sin t}{t}\,dt,&
 \Cin(z)&=\int_0^z\frac{1-\cos t}{t}\,dt,\\
 \Shi(z)&=\int_0^z\frac{\sinh t}{t}\,dt,&
 \Cinh(z)&=\int_0^z\frac{\cosh t-1}{t}\,dt.
\end{aligned}
\end{equation}
Here \(\Cinh(z)=\Chi(z)-\gamma-\Log z\) on a compatible logarithmic
branch; the integral in \eqref{eq:four-entire-integrals} is its entire
normalization.  Termwise comparison gives
\begin{equation}\label{eq:fourth-root-transform}
\begin{aligned}
 \Ein(z)&=\Shi(z)-\Cinh(z),&
 \Ein(-z)&=-\Shi(z)-\Cinh(z),\\
 \Ein(iz)&=\Cin(z)+i\Si(z),&
 \Ein(-iz)&=\Cin(z)-i\Si(z).
\end{aligned}
\end{equation}

\begin{theorem}[Fourth-root orbit]\label{thm:fourth-root-orbit}
For every \(a\in\Qbar^\times\), the six numbers
\begin{equation}\label{eq:six-orbit-values}
 e^a,\quad e^{ia},\quad
 \Si(a),\quad\Cin(a),\quad\Shi(a),\quad\Cinh(a)
\end{equation}
are algebraically independent over \(\Qbar\).  More strongly, the last four
numbers are algebraically independent over \(\Eexp\).
\end{theorem}

\begin{proof}
The four arguments \(a,-a,ia,-ia\) are distinct.  Their \(\Q\)-span has
dimension two, with basis \(a,ia\).  Theorem~\ref{thm:master} therefore
gives transcendence degree six for the corresponding exponential and
\(\Ein\)-coordinates.  The four-by-four transformation
\eqref{eq:fourth-root-transform} is invertible over \(\Qbar\), proving the
first assertion.  The stronger assertion follows directly from
Theorem~\ref{thm:pure-exp-base} and the same invertible transformation.
\end{proof}

\begin{corollary}[A countable independent family]
\label{cor:integer-integral-family}
The countable family
\begin{equation}\label{eq:integer-integral-family}
 \bigl\{\Si(n),\Cin(n),\Shi(n),\Cinh(n):n\ge1\bigr\}
\end{equation}
is algebraically independent over \(\Eexp\).
\end{corollary}

\begin{proof}
It suffices to treat a finite set of positive integers.  The fourth-root
orbits \(\{n,-n,in,-in\}\) are pairwise disjoint as \(n\) varies.  Apply
Theorem~\ref{thm:pure-exp-base} to their union and perform the invertible
transformation \eqref{eq:fourth-root-transform} separately in every block.
\end{proof}

\begin{remark}[Relation to recent sine-integral results]
Delaygue proved that, when
\(a_1^2,b_1^2,\ldots,a_m^2,b_m^2\) are pairwise distinct algebraic
numbers, the integrals
\(\int_{a_j}^{b_j}(\sin t/t)\,dt\) are linearly independent over
\(\Qbar\); see \cite[Proposition~4.4]{Delaygue2025}.  Theorem
\ref{thm:pure-exp-base} and Lemma~\ref{lem:linear-image-ein} upgrade this,
under the same nonzero square-distinctness hypothesis, to algebraic
independence over \(\Eexp\), because
\[
 \int_a^b\frac{\sin t}{t}\,dt
 =\frac{\Ein(ib)-\Ein(-ib)-\Ein(ia)+\Ein(-ia)}{2i}.
\]
This is a strict strengthening of the stated conclusion of Delaygue's
proposition, but a separate priority search in the older
Mahler--Shidlovskii literature is required before presenting the explicit
upgrade as historically new.
\end{remark}

\subsection{The exact first-jet ideal}

The preceding theorem also determines every algebraic relation involving
the first derivatives of the four functions.  This is stronger than a
transcendence-degree count.

\begin{theorem}[Exact first-jet ideal]\label{thm:first-jet-ideal}
Fix \(a\in\Qbar^\times\), and abbreviate
\[
 S=\Si(a),\quad C=\Cin(a),\quad H=\Shi(a),\quad K=\Cinh(a)
\]
and
\[
 S_1=\Si'(a),\quad C_1=\Cin'(a),\quad
 H_1=\Shi'(a),\quad K_1=\Cinh'(a).
\]
The kernel of the evaluation homomorphism
\[
 \Qbar[\mathsf S,\mathsf C,\mathsf H,\mathsf K,
       \mathsf S_1,\mathsf C_1,\mathsf H_1,\mathsf K_1]
 \longrightarrow\C
\]
at \((S,C,H,K,S_1,C_1,H_1,K_1)\) is generated by the two polynomials
\begin{align}
 (a\mathsf S_1)^2+(1-a\mathsf C_1)^2-1,
 \label{eq:trig-jet-relation}\\
 (1+a\mathsf K_1)^2-(a\mathsf H_1)^2-1.
 \label{eq:hyperbolic-jet-relation}
\end{align}
In particular, no nonzero polynomial relation mixes the four zeroth-order
values with the derivative coordinates.
\end{theorem}

\begin{proof}
Put \(x=e^{ia}\) and \(y=e^a\).  Differentiating
\eqref{eq:four-entire-integrals} gives
\begin{equation}\label{eq:jet-torus-coordinates}
\begin{aligned}
 aS_1&=\frac{x-x^{-1}}{2i},&
 1-aC_1&=\frac{x+x^{-1}}2,\\
 aH_1&=\frac{y-y^{-1}}2,&
 1+aK_1&=\frac{y+y^{-1}}2.
\end{aligned}
\end{equation}
Conversely,
\[
 x=1-aC_1+iaS_1,\qquad y=1+aK_1+aH_1,
\]
and the conjugate signs in these formulas give \(x^{-1}\) and \(y^{-1}\).
Thus the derivative coordinate ring is exactly
\(\Qbar[x^{\pm1},y^{\pm1}]\).  Its presentation in the four derivative
variables has precisely the two conic relations
\eqref{eq:trig-jet-relation}--\eqref{eq:hyperbolic-jet-relation}; the
inverse formulas show that there can be no further kernel.  Finally,
Theorem~\ref{thm:fourth-root-orbit} makes \(S,C,H,K\) algebraically
independent over the larger field \(\Eexp\), hence over
\(\Qbar(x,y)\).  Adjoining them is therefore a polynomial extension of the
derivative coordinate ring, and introduces no additional relations.
\end{proof}

\begin{remark}[All positive-order jets]
Repeated differentiation of the four integrands in
\eqref{eq:four-entire-integrals} shows that every derivative of positive
order at \(a\) belongs to \(\Qbar(e^a,e^{ia})\).  Conversely the first
derivatives recover \(e^a,e^{ia}\) by
\eqref{eq:jet-torus-coordinates}.  Hence Theorem
\ref{thm:first-jet-ideal} also says that the four values are algebraically
independent over the field generated by all of their positive-order jets.
\end{remark}

\subsection{A quantitative consequence}

The results above are qualitative.  A recent general theorem supplies
Liouville-type lower bounds, but the exponent is controlled by the
differential transcendence degree, not merely by the number of displayed
integral values.

\begin{proposition}[Quantitative orbit bound]\label{prop:AF-orbit-measure}
Let \(A=\{a_1,\ldots,a_N\}\) be distinct positive real algebraic numbers,
and put
\[
 r=\dim_{\Q}\Span_{\Q}\{a_1,\ldots,a_N\},\qquad
 \tau=4N+2r.
\]
Let \(\boldsymbol\xi_A\) be the \(4N\)-tuple consisting of
\[
 \Si(a_j),\ \Cin(a_j),\ \Shi(a_j),\ \Cinh(a_j)
 \qquad(1\le j\le N).
\]
There exist a constant \(C_2>0\) and positive constants \(C_1(D)>0\),
depending on \(A\) and on the degree bound \(D\), such that every nonzero
polynomial \(P\in\Z[\mathbf X]\) satisfies
\begin{equation}\label{eq:AF-orbit-bound}
 \bigl|P(\boldsymbol\xi_A)\bigr|
 \ge C_1(\deg P)\,
       H(P)^{-C_2(\deg P)^{\tau}}.
\end{equation}
In particular the alternative ``\(P(\boldsymbol\xi_A)=0\)'' in the general
measure theorem never occurs here.
\end{proposition}

\begin{proof}
Consider the \(4N\) E-functions obtained from
\eqref{eq:four-entire-integrals} by replacing their argument by \(a_jz\),
and evaluate them at \(z=1\).  The functional proof of
Theorem~\ref{thm:master}, followed by the blockwise transformation
\eqref{eq:fourth-root-transform}, shows that these functions are
algebraically independent over the exponential differential field generated
by \(e^{\pm a_jz}\) and \(e^{\pm ia_jz}\).  Their derivatives recover that
field, whose torus rank is \(2r\).  Consequently the differential field
generated by the functions and all their derivatives has transcendence
degree \(4N+2r=\tau\) over \(\Qbar(z)\).

Apply the E-function case of
\cite[Theorem~1.3]{AdamczewskiFaverjon2025} at the regular point \(z=1\).
It gives the alternative between vanishing and the lower bound
\eqref{eq:AF-orbit-bound}, with exponent \((\deg P)^\tau\).  Theorem
\ref{thm:pure-exp-base} rules out vanishing for every nonzero \(P\).
\end{proof}

\begin{remark}[Scope of the quantitative import]
Proposition~\ref{prop:AF-orbit-measure} is an application of a general
algebraic-independence measure, not a new measure theorem.  The constants in
\eqref{eq:AF-orbit-bound} are those furnished by
Adamczewski--Faverjon; their paper discusses effectivity separately.  The
same procedure applies to any fixed full-rank linear family supplied by
Lemma~\ref{lem:linear-image-ein}, but its exponent must be computed from the
associated differential field before a numerical value is stated.  No such
quantitative estimate bears on the arithmetic of \(\gamma\) or \(\delta\),
which are connection constants outside the finite-point E-function field.
\end{remark}
